\documentclass[a4paper,12pt]{article}

\usepackage[UTF8]{ctex}
\usepackage{amsmath,amssymb,amsthm,bm}
\usepackage{geometry}
\usepackage{booktabs}
\usepackage{hyperref}
\usepackage{setspace}

\geometry{left=2.5cm,right=2.5cm,top=2.5cm,bottom=2.5cm}
\setstretch{1.25}

\theoremstyle{plain}
\newtheorem{theorem}{定理}[section]
\newtheorem{lemma}[theorem]{引理}
\newtheorem{proposition}[theorem]{命题}
\theoremstyle{definition}
\newtheorem{definition}[theorem]{定义}
\newtheorem{assumption}[theorem]{假设}
\newtheorem{algorithm}[theorem]{算法}
\theoremstyle{remark}
\newtheorem{remark}[theorem]{注}

\title{\textbf{多区域低空交通规划与闭环控制的最小 MFD 模型}}
\author{您的姓名 \\ 您的学院/系, 您的大学}
\date{\today}

\begin{document}

\maketitle

\begin{abstract}
本文面向多区域低空交通（LAAT/UAM）规划与控制，提出一个尽量简洁的宏观闭环模型。模型借鉴多区域 MFD 交通文献的成熟骨架：将城市低空划分为若干近似同质区域，以区域总累积量作为状态，以单峰 MFD 描述区域总产出，以静态路由比例描述跨区域流向，以起飞准入控制调节外部进入流，并以边界放行控制调节区域间转移流。静态层只输出守恒的聚合通量、固定路由比例和区域吞吐需求；动态层只跟踪正吞吐区域的总累积量；控制层显式给出起飞控制器和边界控制器。本文证明：若静态通量守恒、MFD 容量允许目标位于自由流分支、区域间转移矩阵为次随机且谱半径小于一，则活动子系统的目标点是闭环系统的局部指数稳定均衡。该模型不引入显式飞行时延、起降排队、逐目的地组分状态或大规模 MPC，从而给出一个数学自洽且容易仿真实现的多区域低空流量规划与控制主线。

\textbf{关键词：} 低空交通；多区域网络；宏观基本图；路由比例；周界控制；局部稳定性
\end{abstract}

\section{引言}

大规模低空交通网络的核心问题可以概括为两件事：静态上，如何把 OD 服务流量分配到多个空域区域；动态上，如何在扰动下避免某些区域累积量偏离高效运行区间。地面多区域 MFD 文献已经形成了一条成熟建模路线：如果网络能被划分为若干近似同质区域，则每个区域可由总累积量和 MFD 产出函数刻画，周界或边界控制负责调节区域间流量。

Aboudolas 和 Geroliminis 的多水库模型把区域累积量作为状态，用 MFD 产出和边界比例构成守恒方程；Geroliminis、Haddad 和 Ramezani 的两区域 MPC 模型进一步展示了 MFD 控制可用于保护拥堵区域；Yildirimoglu 等与 Sirmatel 和 Geroliminis 将路由指导纳入多区域 MFD 框架；Boufous 等强调单独控制外部进入流不足以保证内部负载分配合理。这些文章的共同点不是复杂的微观航路建模，而是低维区域模型、聚合路由比例和闭环反馈。

因此，本文采用最小建模原则。主模型只保留区域总累积量 $n_i(t)$，不在正文中引入目的地组分状态 $n_i^d(t)$；路径规划只作为静态聚合算法，用来生成起飞计划、跨区域通量和路由比例；动态控制只执行两类低维指令：起飞准入率和边界放行率，而不在线重新决定流向。本文的贡献是把“静态规划--MFD 目标构造--起飞与边界反馈--稳定性判据”压缩为一个完整闭环。

\section{静态路由聚合}

令 $\mathcal R=\{1,\ldots,R\}$ 为低空区域集合，$\mathcal N_i$ 为区域 $i$ 的邻接区域集合。为保持模型简单，本文把邻接关系作为无向支撑集合处理：若两个区域可在任一方向相邻，则二者互为邻接；具体流向由 $T_{ij}$ 或 $\beta_{ij}$ 是否为零决定。静态层从 OD 需求、候选区域路径和拥堵感知路径成本出发，生成聚合通量。

对 OD 对 $m$，令 $q_m$ 为服务流量，$\bar q_m$ 为潜在需求，并要求 $0\le q_m\le \bar q_m<\infty$；令 $\mathcal P_m$ 为非空的候选区域路径集合。给定规划密度估计 $\hat{\bm n}$，路径 $p\in\mathcal P_m$ 的低空飞行时间可取为
\begin{equation}\label{eq:path_time}
    \tau_{mp}(\hat{\bm n})
    =
    \sum_{i\in p}\frac{L_{i,mp}}{v_i(\hat n_i)} ,
\end{equation}
其中 $L_{i,mp}$ 是路径 $p$ 在区域 $i$ 内的航程，$v_i(\cdot)$ 是标定速度函数，并要求在候选规划密度上满足 $v_i(\hat n_i)>0$。路径比例可由 Logit 规则给出：
\begin{equation}\label{eq:logit}
    \alpha_{mp}
    =
    \frac{\exp[-\theta \tau_{mp}(\hat{\bm n})]}
    {\sum_{r\in\mathcal P_m}\exp[-\theta \tau_{mr}(\hat{\bm n})]},
    \qquad \theta>0 .
\end{equation}
服务流量 $q_m$ 可以由弹性需求系统最优、固定需求分配或任意可复现的规划器给出。本文主定理不依赖该规划器的全局最优性，只要求其最终输出可被聚合为守恒的区域通量。

为了使静态输出可以无歧义地进入动态层，本文采用如下聚合约定。把候选路径写成有限区域序列
$p=(i_1,\ldots,i_{\ell_p})$，并令该路径上的服务流为 $f_{mp}=q_m\alpha_{mp}$。若路径从区域 $i_1$ 起飞、在区域 $i_{\ell_p}$ 完成服务，则
\begin{align}
    q_i^g
    &=
    \sum_m\sum_{p\in\mathcal P_m} f_{mp}\mathbf 1\{i_1=i\}, \label{eq:aggregate_q}\\
    T_{ij}
    &=
    \sum_m\sum_{p\in\mathcal P_m} f_{mp}
    \sum_{r=1}^{\ell_p-1}\mathbf 1\{i_r=i,\ i_{r+1}=j\}, \label{eq:aggregate_T}\\
    T_{i0}
    &=
    \sum_m\sum_{p\in\mathcal P_m} f_{mp}\mathbf 1\{i_{\ell_p}=i\}. \label{eq:aggregate_exit}
\end{align}
按式 \eqref{eq:aggregate_q}--\eqref{eq:aggregate_exit} 聚合时，每条路径流在进入某区域、离开某区域或完成服务之间逐段守恒，因此下面的区域守恒式自动成立。若静态规划器直接输出 $q_i^g,T_{ij},T_{i0}$，则必须把这些量校正到同一守恒约束后才能进入动态模型。

令 $q_i^g$ 表示静态层给区域 $i$ 的外部起飞输入，$T_{ij}$ 表示从区域 $i$ 转入邻接区域 $j$ 的静态通量，$T_{i0}$ 表示在区域 $i$ 内完成服务并离开空域动态系统的通量。所有聚合通量均非负，且若 $j\notin\mathcal N_i$ 则取 $T_{ij}=0$。这里的 $0$ 是一个“完成/离网”虚拟节点，不是实际空域区域。静态聚合结果必须满足
\begin{equation}\label{eq:static_conservation}
    q_i^g+\sum_{j\in\mathcal N_i}T_{ji}
    =
    \sum_{j\in\mathcal N_i}T_{ij}+T_{i0},
    \qquad \forall i\in\mathcal R .
\end{equation}

定义区域 $i$ 的静态吞吐需求为
\begin{equation}\label{eq:lambda}
    \Lambda_i
    =
    q_i^g+\sum_{j\in\mathcal N_i}T_{ji}
    =
    \sum_{j\in\mathcal N_i}T_{ij}+T_{i0}.
\end{equation}
若 $\Lambda_i>0$，定义固定路由比例
\begin{equation}\label{eq:beta}
    \beta_{ij}=\frac{T_{ij}}{\Lambda_i},
    \qquad
    \beta_{i0}=\frac{T_{i0}}{\Lambda_i}.
\end{equation}
若 $\Lambda_i=0$，取 $\beta_{ij}=\beta_{i0}=0$。于是对正吞吐区域，
\begin{equation}\label{eq:beta_sum}
    \sum_{j\in\mathcal N_i}\beta_{ij}+\beta_{i0}=1.
\end{equation}
记正吞吐区域集合为
\begin{equation}\label{eq:positive_regions}
    \mathcal R_+=\{i\in\mathcal R:\Lambda_i>0\}.
\end{equation}

\begin{lemma}[零吞吐区域与活动子系统]\label{lem:zero_region}
若 $\Lambda_i=0$，则 $q_i^g=0$、$T_{i0}=0$，且所有进入或离开区域 $i$ 的计划通量均为零。于是按 $\beta_{ij}=\beta_{i0}=0$ 的约定，仍有
\[
    \beta_{ij}\Lambda_i=T_{ij},\qquad \beta_{i0}\Lambda_i=T_{i0}.
\]
特别地，正吞吐区域集合 $\mathcal R_+$ 在固定路由比例下不受零吞吐区域计划流量驱动，也不会向零吞吐区域发送计划流量。
\end{lemma}

\begin{proof}
由式 \eqref{eq:lambda} 和通量非负性，若 $\Lambda_i=0$，则
$q_i^g+\sum_{j\in\mathcal N_i}T_{ji}=0$，所以 $q_i^g=0$ 且所有 $T_{ji}=0$。同理，
$\sum_{j\in\mathcal N_i}T_{ij}+T_{i0}=0$，所以所有 $T_{ij}=0$ 且 $T_{i0}=0$。代入零比例约定即可得到上式。封闭性直接来自进入和离开零吞吐区域的计划通量均为零。
\end{proof}

因此这些零吞吐区域是本模型中的非活动区域。后续稳定性结论只讨论 $\mathcal R_+$ 上的活动子系统；对 $i\notin\mathcal R_+$，仿真时应令 $n_i(0)=0$ 并保持无计划流量，或另设一个独立的安全排空规则，而不把这类区域纳入本文的稳定性定理。

\begin{remark}
拥堵感知路径搜索、Logit 分流和 SO--路径交替迭代都属于静态算法层。若通量由完整路径逐段聚合得到，则式 \eqref{eq:static_conservation} 通常自动成立；若通量由外部规划器直接给出，则必须先投影或校正到守恒约束。若路径算法未收敛，仍可把最后一次输出作为候选路由比例，但必须重新检查式 \eqref{eq:static_conservation}、容量条件和闭环稳定性。本文不把路径固定点收敛作为主模型定理的一部分。
\end{remark}

\section{多区域 MFD 动态}

令 $n_i(t)$ 为区域 $i$ 内航空器总累积量。每个区域的 MFD 产出函数 $G_i(n_i)$ 表示在累积量为 $n_i$ 时区域总流出能力。

\begin{assumption}[一般单峰 MFD]\label{asm:mfd}
对每个区域 $i$，状态满足 $n_i(t)\in\mathbb R_+$，且 $G_i:\mathbb R_+\to\mathbb R_+$。存在临界累积量 $n_i^{\mathrm{cr}}>0$ 和最大产出 $G_i^{\max}>0$，使得 $G_i(0)=0$，$G_i(n)>0$ 对 $n>0$ 成立；$G_i$ 在 $[0,n_i^{\mathrm{cr}}]$ 上连续可微且严格递增，在 $n_i^{\mathrm{cr}}$ 之后严格递减，并满足 $G_i(n_i^{\mathrm{cr}})=G_i^{\max}$。记自由流分支为 $G_{i,s}$，其反函数 $G_{i,s}^{-1}$ 存在。
\end{assumption}

为了区分静态计划和动态执行，令 $r_i(t)$ 表示区域 $i$ 的实际起飞准入率，令 $z_i(t)$ 表示区域 $i$ 的实际总边界放行率。采用静态路由比例 $\beta_{ij}$ 分配边界流向：
\begin{equation}\label{eq:boundary_flow}
    z_{ij}(t)=\beta_{ij}z_i(t),
    \qquad
    z_{i0}(t)=\beta_{i0}z_i(t).
\end{equation}
其中 $z_{ij}$ 是从区域 $i$ 进入邻接区域 $j$ 的执行边界流，$z_{i0}$ 是完成服务并离开系统的执行流。这里作一个显式建模假设：$z_i$ 是区域总移除/转移调度率，既包含跨区转移，也包含完成服务后的离网移除；二者均按静态比例执行。于是多区域总量动态为
\begin{equation}\label{eq:dynamics}
    \dot n_i(t)
    =
    r_i(t)
    +
    \sum_{j\in\mathcal N_i}\beta_{ji}z_j(t)
    -
    z_i(t).
\end{equation}
该式是本文的核心被控对象。它只描述区域总累积量，不追踪每个目的地组分；完成流由 $z_{i0}(t)$ 显式表示。由于 $r_i(t)\ge0$、$z_j(t)\ge0$ 且 $G_i(0)=0$，当 $n_i=0$ 时有 $z_i=0$，故非负正交域 $\mathbb R_+^R$ 对闭环动态正不变。控制律只对活动区域 $i\in\mathcal R_+$ 定义。若需要书写全区域仿真，可对 $i\notin\mathcal R_+$ 取 $r_i(t)=z_i(t)=0$ 并令 $n_i(0)=0$，此时零吞吐区域保持在零状态；若非活动区域存在初始航空器，则必须另设排空规则，不能由本文稳定性定理覆盖。由引理 \ref{lem:zero_region}，对活动子系统 $i\in\mathcal R_+$ 而言，来自 $j\notin\mathcal R_+$ 的计划转移项为零，因此后续稳定性分析可把求和等价地限制在 $\mathcal R_+$ 内。

\section{目标与可行性}

选取边界放行裕度 $\bar u_i\in(0,1)$。该裕度使目标边界放行比例不位于饱和边界，从而在目标附近存在未饱和反馈邻域。对正吞吐区域，若
\begin{equation}\label{eq:capacity}
    0<\Lambda_i<\bar u_iG_i^{\max},
\end{equation}
则定义自由流目标累积量
\begin{equation}\label{eq:target}
    \bar n_i
    =
    G_{i,s}^{-1}\left(\frac{\Lambda_i}{\bar u_i}\right).
\end{equation}
若 $\Lambda_i=0$，取 $\bar n_i=0$，并把该区域作为非活动区域处理；它不参与下面的活动子系统稳定性分析。
此外，令 $\bar q_i^{\mathrm{phys}}$ 表示区域 $i$ 的物理起飞容量，$\bar q_i^{\mathrm{pot}}$ 表示当前可被接受的潜在起飞需求上限，并定义可执行起飞准入上限
\begin{equation}\label{eq:takeoff_capacity}
    \bar q_i^a=\min\{\bar q_i^{\mathrm{phys}},\bar q_i^{\mathrm{pot}}\}.
\end{equation}
静态计划必须满足 $0\le q_i^g\le \bar q_i^a$。若要在起飞端使用严格的局部反馈增益，则要求 $0<q_i^g<\bar q_i^a$；若没有起飞上调或下调余量，则把该区域的起飞反馈增益设为零，只执行静态计划起飞率。

若实际边界存在逐界面容量 $C_{ij}\in(0,\infty]$，其中 $j\in\mathcal N_i\cup\{0\}$，则静态边界计划还需满足严格局部裕度
\begin{equation}\label{eq:interface_capacity}
    T_{ij}<C_{ij},\qquad \forall i\in\mathcal R_+,\quad j\in\mathcal N_i\cup\{0\}.
\end{equation}
当某一界面不作为容量约束建模时，可取 $C_{ij}=\infty$。由于局部未饱和时 $z_{ij}(t)=\beta_{ij}(\Lambda_i+\kappa_i e_i(t))$ 且 $T_{ij}=\beta_{ij}\Lambda_i$，式 \eqref{eq:interface_capacity} 保证存在目标邻域使逐界面执行流仍满足 $z_{ij}(t)\le C_{ij}$。

\begin{proposition}[目标可构造性]\label{prop:target}
在假设 \ref{asm:mfd} 和容量条件 \eqref{eq:capacity} 下，每个正吞吐区域存在唯一自由流目标 $\bar n_i\in(0,n_i^{\mathrm{cr}})$，并满足
\begin{equation}\label{eq:target_balance}
    \bar u_iG_i(\bar n_i)=\Lambda_i.
\end{equation}
\end{proposition}

\begin{proof}
由假设 \ref{asm:mfd}，$G_{i,s}$ 在 $[0,n_i^{\mathrm{cr}}]$ 上严格递增，且值域为 $[0,G_i^{\max}]$。容量条件给出 $\Lambda_i/\bar u_i\in(0,G_i^{\max})$，因此自由流反函数给出唯一 $\bar n_i=G_{i,s}^{-1}(\Lambda_i/\bar u_i)$，并且 $\bar n_i$ 位于自由流分支内部。
\end{proof}

\begin{proposition}[静态--动态桥接]\label{prop:bridge}
若静态通量满足守恒式 \eqref{eq:static_conservation}，路由比例由式 \eqref{eq:beta} 给出，且正吞吐区域的目标由式 \eqref{eq:target} 构造，则在
\[
    \bar r_i=q_i^g,\qquad
    \bar z_i=\Lambda_i,\qquad
    \bar z_{ij}=T_{ij},\qquad
    \bar z_{i0}=T_{i0}
\]
下，$\bar{\bm n}$ 是活动子系统 \eqref{eq:dynamics} 的稳态。若同时令所有非活动区域满足 $n_i(0)=\bar n_i=0$ 且无计划流量，则它们不影响该稳态。
\end{proposition}

\begin{proof}
固定任意 $i\in\mathcal R_+$。由引理 \ref{lem:zero_region}，$j\notin\mathcal R_+$ 的转移项不进入活动子系统；对 $j\in\mathcal R_+$，目标点满足 $\bar z_j=\Lambda_j$，且起飞准入为 $\bar r_i=q_i^g$。代入式 \eqref{eq:dynamics} 得
\begin{equation}
    \dot n_i
    =
    q_i^g+\sum_{j\in\mathcal N_i\cap\mathcal R_+}\beta_{ji}\Lambda_j-\Lambda_i.
\end{equation}
由 $\beta_{ji}\Lambda_j=T_{ji}$、引理 \ref{lem:zero_region} 对零吞吐区域的约定以及定义 \eqref{eq:lambda}，右端等于 $q_i^g+\sum_jT_{ji}-\Lambda_i=0$。故目标点为活动子系统的稳态。
\end{proof}

\section{反馈控制与稳定性}

定义误差 $e_i(t)=n_i(t)-\bar n_i$。本文采用两个执行控制器。

第一，起飞准入控制器调节外部进入流：
\begin{equation}\label{eq:takeoff_controller}
    r_i(t)
    =
    \operatorname{sat}_{[0,\bar q_i^a]}
    \left(q_i^g-\eta_i e_i(t)\right),
    \qquad
    \eta_i\ge0 .
\end{equation}
其中 $\bar q_i^a$ 是区域 $i$ 的最大可执行起飞准入率。若 $\eta_i>0$，需要 $q_i^g$ 位于区间 $(0,\bar q_i^a)$ 内部，从而目标附近起飞控制不饱和；若该内点条件不满足，则取 $\eta_i=0$，即只执行静态计划起飞率。未被 $r_i(t)$ 接受的潜在需求停留在空域外，本文不把地面或起降场排队纳入状态方程。

第二，边界控制器调节区域总边界放行率。对每个正吞吐区域，采用带分母保护的比例放行律：
\begin{equation}\label{eq:boundary_controller}
    u_i(t)
    =
    \operatorname{sat}_{[0,1]}
    \left(
    \frac{\Lambda_i+\kappa_i e_i(t)}
    {\max\{G_i(n_i(t)),\epsilon_i\}}
    \right),
    \qquad
    \kappa_i>0,\quad 0<\epsilon_i<G_i(\bar n_i).
\end{equation}
并令
\begin{equation}\label{eq:z_controller}
    z_i(t)=u_i(t)G_i(n_i(t)),
    \qquad
    z_{ij}(t)=\beta_{ij}z_i(t),
    \qquad
    z_{i0}(t)=\beta_{i0}z_i(t).
\end{equation}
因此，$u_i$ 是区域汇总边界放行比例，$z_i$ 是实际总放行流，$z_{ij}$ 是可下发到边界 $i\to j$ 的执行命令。

当轨迹位于目标附近、控制未饱和且 $G_i(n_i)>\epsilon_i$ 时，
\begin{equation}\label{eq:controlled_outflow}
    r_i(t)=q_i^g-\eta_i e_i(t),
    \qquad
    z_i(t)=\Lambda_i+\kappa_i e_i(t).
\end{equation}

在目标点，
\[
    \frac{\Lambda_i}{G_i(\bar n_i)}=\bar u_i\in(0,1),
    \qquad G_i(\bar n_i)>\epsilon_i .
\]
令
\[
    \phi_i(n_i)=\frac{\Lambda_i+\kappa_i(n_i-\bar n_i)}{G_i(n_i)} .
\]
由于 $\phi_i(\bar n_i)=\bar u_i$，且 $G_i$ 在目标附近连续可微并满足 $G_i(\bar n_i)>\epsilon_i$，存在 $\delta_i>0$，使得当 $|n_i-\bar n_i|<\delta_i$ 时，
$G_i(n_i)>\epsilon_i$ 且 $0<\phi_i(n_i)<1$。若 $\eta_i>0$，由 $0<q_i^g<\bar q_i^a$ 也可取足够小的邻域使 $0<q_i^g-\eta_i e_i<\bar q_i^a$；若 $\eta_i=0$，起飞准入恒等于 $q_i^g$。对有限个活动区域取这些邻域的交集，即得到一个共同的未饱和邻域。也就是说，式 \eqref{eq:controlled_outflow} 是一个局部等式；若轨迹离开该邻域，饱和控制仍然有定义，但本文的稳定性定理不声称全局稳定。

令 $B=[\beta_{ij}]_{i,j\in\mathcal R_+}$ 为只包含正吞吐区域间转移比例的矩阵，不包含完成/离网比例 $\beta_{i0}$。由于 $\beta_{ij}\ge0$，且由引理 \ref{lem:zero_region} 可把非活动区域转移项置零，所以对 $i\in\mathcal R_+$ 有 $\sum_{j\in\mathcal R_+}\beta_{ij}+\beta_{i0}=1$。因此 $B$ 是非负次随机矩阵；若区域 $i$ 有正完成流，则第 $i$ 行行和严格小于 $1$。令
\[
    K=\operatorname{diag}(\kappa_i)_{i\in\mathcal R_+},
    \qquad
    H=\operatorname{diag}(\eta_i)_{i\in\mathcal R_+},
\]
误差向量 $\bm e$ 也限制在 $\mathcal R_+$ 上。

\begin{theorem}[闭环局部指数稳定]\label{thm:stability}
考虑所有正吞吐区域 $\mathcal R_+$ 组成的活动子系统。假设：
\begin{enumerate}
    \item 静态通量非负且满足守恒式 \eqref{eq:static_conservation}；
    \item 对每个 $i\in\mathcal R_+$，容量条件 $0<\Lambda_i<\bar u_iG_i^{\max}$ 成立，起飞计划满足 $0\le q_i^g\le \bar q_i^a$，且若 $\eta_i>0$ 则 $0<q_i^g<\bar q_i^a$；
    \item 对所有受容量约束的执行界面，逐界面裕度 \eqref{eq:interface_capacity} 成立；
    \item 非活动区域满足零初值并保持无计划流量，或由独立排空规则处理而不进入活动子系统；
    \item 控制参数满足式 \eqref{eq:takeoff_controller}--\eqref{eq:z_controller} 的条件；
    \item 谱半径条件成立：
    \begin{equation}\label{eq:rho}
        \rho(B)<1 .
    \end{equation}
\end{enumerate}
则存在一个目标邻域，使目标均衡点 $\bar{\bm n}$ 对活动子系统上的闭环系统 \eqref{eq:dynamics}, \eqref{eq:takeoff_controller}--\eqref{eq:z_controller} 局部指数稳定。
\end{theorem}

\begin{proof}
在未饱和邻域内，由式 \eqref{eq:controlled_outflow} 和稳态桥接关系可得
\begin{equation}
    \dot e_i
    =
    \sum_{j\in\mathcal N_i\cap\mathcal R_+}\beta_{ji}\kappa_j e_j-\kappa_i e_i-\eta_i e_i .
\end{equation}
写成向量形式为
\begin{equation}\label{eq:error_system}
    \dot{\bm e}
    =
    \left((B^T-I)K-H\right)\bm e .
\end{equation}
因为 $B\ge0$ 且 $\rho(B)<1$，$I-B^T$ 是非奇异 $M$-矩阵。为说明右乘正对角矩阵后该性质仍保持，取 $x>0$ 使 $(I-B^T)x>0$，令 $y=K^{-1}x>0$，则 $(I-B^T)Ky=(I-B^T)x>0$；又 $(I-B^T)K$ 仍是 $Z$-矩阵。由 $Z$-矩阵的正向量判据，若某个 $Z$-矩阵 $M$ 存在正向量 $y>0$ 使 $My>0$，则 $M$ 为非奇异 $M$-矩阵，因此 $(I-B^T)K$ 是非奇异 $M$-矩阵。再加上非负对角矩阵 $H$ 后，$((I-B^T)K+H)y>0$ 且 $(I-B^T)K+H$ 仍为 $Z$-矩阵，所以它仍为非奇异 $M$-矩阵，其全部特征值实部为正。因此
\[
    A=-\bigl((I-B^T)K+H\bigr)=(B^T-I)K-H
\]
为 Hurwitz。

取任意正定矩阵 $Q=I$，由 Hurwitz 性存在唯一正定矩阵 $P$ 满足 Lyapunov 方程
\begin{equation}
    A^TP+PA=-Q .
\end{equation}
令 $V(\bm e)=\bm e^TP\bm e$，则在线性误差系统 \eqref{eq:error_system} 上有 $\dot V=-\bm e^TQ\bm e<0$。因此误差系统指数稳定。在上述未饱和邻域内，闭环误差动力学精确等于式 \eqref{eq:error_system}，故目标点局部指数稳定。
更具体地，可取一个足够小的 Lyapunov 子水平集 $\{\bm e:V(\bm e)\le c\}$ 使其完全包含在共同未饱和邻域内。由于 $\dot V<0$，该子水平集正不变，轨迹不会离开未饱和邻域，所以局部指数稳定结论与原闭环系统一致，而不是只对形式线性化成立。
\end{proof}

\begin{remark}
条件 $\rho(B)<1$ 的物理含义是区域间转移最终存在“泄漏”，即每个闭合通信类都能把一部分流量完成服务并离开多区域空域系统。它不是由次随机性自动推出的条件，而是需要显式检查的条件。若存在一个封闭区域环路，且环路内所有区域都没有完成流或离网流，则 $\rho(B)$ 可能达到 $1$，此时仅靠固定路由比例和总量放行控制无法保证所有区域误差衰减。
\end{remark}

\section{算法与数值验证}

\begin{algorithm}[最小闭环验证流程]\label{alg:verify}
给定区域划分、候选路径、需求和 MFD 标定对象，按以下步骤完成规划与控制验证：
\begin{enumerate}
    \item 用式 \eqref{eq:path_time}--\eqref{eq:logit} 或其他静态规划器生成 OD 服务流量与路径比例，并保证每条候选路径是有限区域序列。
    \item 按式 \eqref{eq:aggregate_q}--\eqref{eq:aggregate_exit} 聚合得到 $q_i^g$、$T_{ij}$ 和 $T_{i0}$；若使用外部规划器直接输出通量，则先校正到静态守恒式 \eqref{eq:static_conservation}。
    \item 计算 $\Lambda_i$、$\beta_{ij}$ 和 $\beta_{i0}$，检查路由比例归一化 \eqref{eq:beta_sum}，并识别活动区域集合 $\mathcal R_+$。对非活动区域，确认无计划流量且仿真初值为零，或把它们交给独立排空规则。
    \item 对每个 $i\in\mathcal R_+$ 选择边界放行裕度 $\bar u_i$，检查 $\Lambda_i<\bar u_iG_i^{\max}$；同时检查起飞计划满足 $0\le q_i^g\le \bar q_i^a$，逐界面计划满足 $T_{ij}<C_{ij}$，$j\in\mathcal N_i\cup\{0\}$。若检查失败，则先降低需求、改变路由或重新划分区域；若通过，则计算目标 $\bar n_i$。
    \item 选择起飞反馈增益 $\eta_i$、边界反馈增益 $\kappa_i$ 和正则化参数 $\epsilon_i<G_i(\bar n_i)$。若 $q_i^g$ 不在 $(0,\bar q_i^a)$ 内部，则取 $\eta_i=0$。
    \item 在 $\mathcal R_+$ 上构造 $B$，检查 $\rho(B)<1$。若谱半径条件失败，则需要增加完成/离网通量或重新设计静态路由。
    \item 在连续时间仿真或 ODE 求解器中按式 \eqref{eq:takeoff_controller} 生成起飞准入命令 $r_i(t)$，按式 \eqref{eq:boundary_controller}--\eqref{eq:z_controller} 生成总边界放行命令 $z_i(t)$ 和逐边界命令 $z_{ij}(t)$。若工程实现采用采样保持控制，则额外检查采样周期 $\Delta t$ 下的离散化闭环矩阵仍满足谱半径小于 $1$。
    \item 只在目标附近施加小扰动来验证局部结论，记录 $\|\bm n(t)-\bar{\bm n}\|$、实际起飞准入率、完成流、边界控制饱和频率和区域最大累积量；若出现频繁饱和，应报告为超出本文局部稳定性条件。
\end{enumerate}
\end{algorithm}

建议数值部分只保留三组实验。第一，闭环稳定实验：从目标附近不同初值出发，验证误差指数衰减并对比谱横坐标，同时记录 $r_i(t)$ 与 $z_{ij}(t)$ 的执行轨迹。第二，路由策略实验：比较固定最短路、拥堵感知 Logit 路由和人工均衡路由对 $\Lambda_i$、$\rho(B)$ 和最大区域累积量的影响。第三，鲁棒与饱和实验：扰动 MFD 容量、需求水平、起飞准入上限和控制裕度，观察何时违反容量条件或进入饱和区。

\section{模型边界}

本文有意不处理以下扩展。第一，显式飞行时延未进入动态方程；若区域间飞行时间与控制采样周期同量级，应引入延迟状态或 MPC。第二，未被起飞准入控制器接受的潜在需求停留在空域外，起降场或地面排队稳定性不作为本文闭环稳定性的一部分。第三，目的地组分状态可用于更细的路由一致性分析，但会显著增加稳定性证明复杂度，本文把它留作附录或后续工作。第四，在线重路由会把固定 $B$ 的系统变成切换系统，需要新的稳定性条件。第五，当区域已经处于拥堵分支并满足 $G_i(n_i)<\Lambda_i$ 时，提高放行率也可能无法恢复，需要需求削减或临时禁入。第六，本文的稳定性结论是连续时间反馈下的目标附近局部结论；若初值远离目标并触发饱和，控制律仍可执行，但不能仅由定理 \ref{thm:stability} 保证恢复。若采用采样保持实现，需要额外检查采样后的离散闭环稳定性。第七，零吞吐区域若存在非零初始累积量，不会被本文的活动子系统控制律自动排空，必须作为单独安全处理模块。第八，逐界面容量只通过局部裕度条件 \eqref{eq:interface_capacity} 覆盖；若容量随时间变化或与微观冲突约束耦合，则需要额外安全层。

\section{结论}

本文给出一个最小的多区域低空交通 MFD 闭环模型。静态层提供守恒通量和路由比例，MFD 自由流分支给出目标累积量，起飞准入控制器执行外部进入流调节，边界控制器执行区域总放行和逐边界流量分解，谱半径条件 $\rho(B)<1$ 给出直接可算的局部稳定性判据。模型的闭环对象明确限制在正吞吐活动区域上，非活动区域只在零初值或独立排空规则下进入仿真，因此静态规划、目标构造、起飞与边界反馈、稳定性判据之间形成一个简单闭环。与更复杂的逐路径、逐目的地、逐边界模型相比，该模型更接近多区域 MFD 文献的主流简洁结构，也更适合作为后续仿真和论文写作的核心理论框架。

\end{document}
